\section{Uso básico de Canvas}

Comenzamos observando el elemento \taghtml{\textless canvas\textgreater}. Al final de esta sección, sabrás como configurar el entorno 2D de \code{canvas} y habrás dibujado el primer ejemplo en tu navegador.

\subsection{El elemento canvas}

\begin{lstlisting}[language=HTML, style=mystyle2]
  <canvas id="tutorial" width="150" height="150"></canvas>
\end{lstlisting}

A primera vista, un elemento \taghtml{\textless canvas\textgreater} es parecido al elemento \taghtml{\textless img\textgreater}, con la diferencia que este no tiene los atributos \code{src} y \code{alt}. El elemento \taghtml{\textless canvas\textgreater} tiene solo dos atributos - \code{width} y \code{height}. Ambos son opcionales y pueden ser definidos usando propiedades del \lineCode{DOM}. Cuando los atributos ancho y alto no están especificados, el lienzo se inicializara con 300 pixels de ancho y 150 pixels de alto. El elemento puede ser arbitrariamente re-dimensionado por \textcss{CSS}, pero durante el rende rizado la imagen es escalada para ajustarse al tamaño de su layout. Si el tamaño del \textcss{CSS} no respeta el ratio del canvas inicial, este aparecerá distorsionado.

\vspace{0.5cm} % separación vertical
\begin{tcolorbox}
	[colback=red!5!white,colframe=cyan,fonttitle=\bfseries,title={\faLightbulbO\, Nota:}]
	Si su renderizado se ve distorsionado, pruebe especificar los atributos \code{width} y \code{height} \\ explícitamente en los atributos del \taghtml{\textless canvas\textgreater} , y no usando \textcss{CSS}.
\end{tcolorbox}
\vspace{0.5cm} % separación vertical

El atributo \code{id} no está especificado para el elemento \taghtml{\textless canvas\textgreater} pero es uno de los atributos globales de HTML el cual puede ser aplicado a cualquier elemento HTML (como \code{class} por ejemplo). Siempre es buena idea proporcionar un \code{id} porque esto hace más fácil identificarlo en un \code{script}.

El elemento \taghtml{\textless canvas\textgreater} puede ser estilizado como a cualquier imagen normal (\code{margin}, \code{border}, \code{background}, etc). Estas reglas, sin embargo, no afectan a lo dibujado sobre el canvas. Cuando no tenemos reglas de estilo aplicadas al \code{canvas}, este será completamente transparente.


\subsubsection{Contenido alternativo}

El elemento \taghtml{\textless canvas\textgreater} se diferencia de un tag \taghtml{\textless img\textgreater} en que, como los elementos \taghtml{\textless video\textgreater}, \taghtml{\textless audio\textgreater} o \taghtml{\textless picture\textgreater}, es fácil definir contenido alternativo (fallback content) para mostrarse en navegadores viejos que no soporten el elemento \taghtml{\textless canvas\textgreater}, como versiones de Internet Explorer previas a la versión 9 o navegadores de texto. Siempre debes proporcionar contenido alternativo para mostrar en estos navegadores.

Proporcionar contenido alternativo es muy explicito: solo debemos insertar el contenido alterno dentro del elemento \taghtml{\textless canvas\textgreater}. Los navegadores que no soporten \taghtml{\textless canvas\textgreater} ignoraran el contenedor y mostrarán el contenido indicado dentro de este. Y los que soporten \taghtml{\textless canvas\textgreater} ignorarán el contenido en su interior (de las etiquetas), y mostrarán el \code{canvas} normalmente.

Por ejemplo, podremos proporcionar un texto descriptivo del contenido del \code{canvas} o proveer una imagen estática del contenido rederizado. Nos podría quedar algo así:

\vspace{0.5cm} % separación vertical
\begin{lstlisting}[language=HTML, style=mystyle2]
  <canvas id="stockGraph" width="150" height="150">
    current stock price: $3.15 +0.15
  </canvas>

  <canvas id="clock" width="150" height="150">
    <img src="images/clock.png" width="150" height="150" alt=""/>
  </canvas>
\end{lstlisting}

\subsection{Etiqueta canvas requerida}

De manera distinta al elemento \taghtml{\textless img\textgreater},, el elemento \taghtml{\textless canvas\textgreater} requiere cerrar la etiqueta con (\taghtml{\textless /canvas\textgreater}).

\vspace{0.5cm} % separación vertical
\begin{tcolorbox}
	[colback=red!5!white,colframe=cyan,fonttitle=\bfseries,title={\faLightbulbO\, Nota:}]

	Aunque las versiones anteriores del navegador \textblue{Safari} de Apple no requeria el cierre de la etiqueta, la especificacion indica que es necesaria, asi que tu deberias incluir esta para asegurarte la compatibilidad. Aquellas versiones de Safari (anteriores versiones a 2.0) renderizaran el contenido de regreso agregandolo al canvas mismo a no ser que utilice trucos de \textcss{CSS} para enmascararlo. Afortunadamente, los usuarios de aquellas versiones de Safari son raros hoy en dia.
\end{tcolorbox}
\vspace{0.5cm} % separación vertical

Si el contenido alternativo no se necesita, un simple \lineCode{\textless canvas id="foo" \ldots \textgreater\textless /canvas\textgreater} es completamente compatible con todos los navegadores que soportan canvas.

\newpage % saltar a la pagina siguiente
\subsection{El contexto de renderización}

\taghtml{\textless canvas\textgreater} crea un lienzo de dibujo fijado que expone uno o mas contextos renderizados, los cuales son usados para crear y manipular el contenido mostrado. Nos enfocaremos en renderizacion de contextos 2D. Otros contextos deberán proveer diferentes tipos de renderizaciones; por ejemplo, \textblue{WebGL} usa un 3D contexto (``experimental-webgl'') basado sobre \textblue{OpenGL ES}.

El \code{canvas} esta inicialmente en blanco. Para mostrar alguna cosa, un \code{script} primero necesita acceder al contexto a renderizar y dibujar sobre este. El elemento \taghtml{\textless canvas\textgreater} tiene un method llamado \code{getContext()}, usado para obtener el contexto a renderizar y sus funciones de dibujo. \code{getContext()} toma un parámetro, el tipo de contexto. Para graficos 2D, como los que cubre este tutorial, su especificacion es "2d".

\vspace{0.5cm} % separación vertical
\begin{lstlisting}[language=TypeScript, style=mystyle]
  const canvas = document.getElementById("tutorial");
  const ctx = canvas.getContext("2d");
\end{lstlisting}
\vspace{0.5cm} % separación vertical

La primera linea regresa el nodo \lineCode{DOM} para el elemento \taghtml{\textless canvas\textgreater} llamando al método \\ \code{document.getElementById()}. Una vez tu tienes el elemento nodo, tu puedes acceder al contexto de dibujo usando su método \code{getContext()}.

\subsection{Comprobando soporte}

El contenido de regreso que es mostrado en navegadores los cuales no soportan \taghtml{\textless canvas\textgreater}. Para los \code{Scripts} puede también comprobarse su soporte desde la programación por un simple test para la presencia del método \code{getContext()}. Con un trozo de código parecido al que viene debajo:

\vspace{0.5cm} % separación vertical
\begin{lstlisting}[language=TypeScript, style=mystyle]
  const canvas = document.getElementById("tutorial");

  if (canvas.getContext) {
    let ctx = canvas.getContext("2d");
    // drawing code here
  } else {
    // canvas-unsupported code here
  }
\end{lstlisting}
\vspace{0.5cm} % separación vertical

\newpage % saltar a la siguiente pagina

\subsection{Un esqueleto de plantilla}

Aquí esta una plantilla minimalista, la cual usaremos como punto de partida para posteriores ejemplos.

\vspace{0.5cm} % separación vertical
\begin{lstlisting}[language=HTML, style=mystyle2]
  <html>
    <head>
      <title>Canvas tutorial</title>
      <script type="text/javascript">
        function draw() {
          const canvas = document.getElementById("tutorial");
          if (canvas.getContext) {
            let ctx = canvas.getContext("2d");
          }
        }
      </script>
      <style type="text/css">
        canvas {
          border: 1px solid black;
        }
      </style>
    </head>
    <body onload="draw();">
      <canvas id="tutorial" width="150" height="150"></canvas>
    </body>
  </html>
\end{lstlisting}
\vspace{0.5cm} % separación vertical

El \code{script} incluye una función llamada \code{draw()}, la cual es ejecutada una vez finalizada la carga de la pagina; este esta hecho usando el evento \code{load} del documento. Esta función, o una parecida, podría también ser llamada usando \code{window.setTimeout()} , \code{window.setInterval()} , o cualquier otro manejador de evento, a lo largo de que la pagina esta siendo cargada la primera vez.

Así es como se ve el elemento \taghtml{\textless canvas\textgreater} en la pagina.

\vspace{0.5cm} % separación vertical
\begin{center}
	\begin{tikzpicture}
		\draw (0,0) rectangle (5,5); % Cuadrado de esquina (0,0) a (2,2)
	\end{tikzpicture}
\end{center}
\vspace{0.5cm} % separación vertical

\newpage % saltar a la siguiente pagina
\subsection{Ejemplo}

Para comenzar, daremos un vistazo a un simple ejemplo que dibuja dos rectángulos que se intersectan, uno de los cuales tiene alpha transparencia.

\vspace{0.5cm} % separación vertical
\begin{lstlisting}[language=HTML, style=mystyle2]
  <html>
    <head>
      <script type="application/javascript">
        function draw() {
          const canvas = document.getElementById("canvas");
          if (canvas.getContext) {
            let ctx = canvas.getContext("2d");

            ctx.fillStyle = "rgb(200,0,0)";
            ctx.fillRect(10, 10, 55, 50);

            ctx.fillStyle = "rgba(0, 0, 200, 0.5)";
            ctx.fillRect(30, 30, 55, 50);
          }
        }
      </script>
    </head>
    <body onload="draw();">
      <canvas id="canvas" width="150" height="150"></canvas>
    </body>
  </html>
\end{lstlisting}
\vspace{0.5cm} % separación vertical

Visualmente quedaría así:

\vspace{0.5cm} % separación vertical
\begin{center}
	\begin{tikzpicture}
		% Rectángulo rojo
		\fill[red] (0, 0) rectangle (5, 5);

		% Rectángulo azul semi-transparente
		\fill[blue, opacity=0.5] (2.5, 2.5) rectangle (7.5, -2.5);
	\end{tikzpicture}
\end{center}
